\begin{textsample}{POS Dim 4 – 2021\_c5 – Score 1.00 – 2021\_c5/enade\_2021\_mat\_6642fe63dd.txt}  \label{ex:f4_pos_215}
CAPÍTULO 5 DISTRIBUIÇÃO DOS CONCEITOS 5.1 PANORAMA NACIONAL DA DISTRIBUIÇÃO DOS CONCEITOS O Conceito Enade é um indicador de qualidade que avalia os cursos por intermédio dos desempenhos dos estudantes a partir dos resultados do Enade.
Os valores possíveis vão de 1 ( pior situação ) a 5 ( melhor situação ).
Os cursos com apenas um ou sem nenhum concluinte participante não obtêm o Conceito Enade, ficando Sem Conceito ( SC ).
Na Tabela 5.1, são apresentadas a quantidade e a distribuição de cursos de Matemática ( Licenciatura ) participantes do Enade/2021, por faixa de conceito e Grande Região.
Enfatiza -se, mais uma vez, que a diferença entre os cursos tabulados neste capítulo e no capítulo 2 corresponde aos cursos que não foram avaliados, ou seja, àqueles, em princípio, sem alunos concluintes que tenham sido inscritos para a prova.
Observando -se os dados da Tabela 5.1, nota -se que, dos 435 cursos participantes, 153 ( 35,2 \% ) classificaram -se com conceito 2, o valor modal.
Esse foi também o conceito modal na maioria das regiões : Norte ( 64,7 \% ), Nordeste ( 43,2 \% ) e Centro-Oeste ( 44,9 \% ).
O conceito 3 foi o segundo mais frequente em nível nacional ( 34,0 \%, correspondendo a 148 cursos ), tendo sido o conceito modal da região Sul ( 38,5 \% ).
O conceito 4 foi o terceiro mais frequente ( 19,1 \%, correspondendo a 83 cursos ), tendo sido o conceito modal da região Sudeste ( 31,7 \% ).
Houve, ainda, 17 cursos ( 3,9 \% ) que receberam o conceito 5, e 19 cursos ( 4,4 \% ) que receberam o conceito 1.
Dos 435 cursos de Matemática ( Licenciatura ), 15 ( 3,4 \% ) ficaram Sem Conceito ( SC ).
A região Norte participou com 34 cursos, o correspondente a 7,8 \% do total nacional.
Desses, 22 cursos receberam o conceito 2, o que equivale a 64,7 \% do total regional, além de ter sido o conceito modal na região, como já comentado.
A quatro cursos ( 11,8 \% ) atribuiu -se No Gráfico 5.1, é apresentada a distribuição cumulativa do Conceito Enade, segundo a Grande Região.
Em um gráfico de distribuição cumulativa, em cada valor de abscissa ( eixo x ), o valor apresentado no eixo das ordenadas ( eixo y ) é o valor acumulado da variável ( Conceito Enade ).
Observando -se a linha vertical pontilhada, que ajuda a visualizar o percentual de conceitos acumulados até o valor 3, ou seja, os conceitos 1, 2 e 3, pode -se notar que, por exemplo, a região Norte ( linha azul ) apresentou 100 \% dos cursos com o conceito 3 ou abaixo desse valor.
A região Sul ( linha rosa ), por outro lado, apresentou 61,8 \% dos cursos na mesma situação : com o conceito 3 ou abaixo.
Destaca -se que, em algumas regiões, o conceito 5 pode não ter sido alcançado por nenhum curso.
Foi o que ocorreu na região Norte ( linha azul ), na qual todos os cursos obtiveram conceitos menores ou iguais a 3, tendo a poligonal alcançado 100 \% no conceito 3, e na região Centro-Oeste ( linha laranja ), na o conceito 1, e a seis cursos ( 17,6 \% ), o conceito 3.
Dois cursos ficaram Sem Conceito ( SC ), e nenhum curso recebeu os conceitos 4, ou 5 nessa região.
A região Nordeste participou com 132 cursos, o equivalente a 30,3 \% do total nacional.
Nessa região, 57 cursos ( 43,2 \% em termos regionais ) obtiveram conceito 2, o conceito modal para a região, como já comentado.
Os conceitos 1 e 3 foram atribuídos, respectivamente, a oito e a 55 cursos ( 6,1 \% e 41,7 \%, respectivamente ).
Ainda nessa região, oito cursos ( 6,1 \% ) foram avaliados com o conceito 4, e um curso ( 0,8 \% ) foi avaliado com o conceito 5.
Três cursos ficaram Sem Conceito ( SC ).
Dos 142 cursos participantes da região Sudeste ( 32,6 \% do total nacional ), o conceito 4 foi o conceito modal, como já anteriormente assinalado, tendo sido obtido por 45 cursos ( 31,7 \% ).
O conceito 1 foi atribuído a dois cursos ( 1,4 \% ), e o conceito 2, a 36 cursos ( 25,4 \% ).
Receberam os conceitos 3 e 5, respectivamente, 43 cursos ( 30,3 \% ) e dez cursos ( 7,0 \% ).
Nessa região, seis cursos ficaram Sem Conceito ( SC ).
Os 78 cursos da região Sul corresponderam a 17,9 \% do total nacional.
Como já indicado, houve predominância do conceito 3, atribuído a 38,5 \% dos cursos da região ( conceito modal ), ou seja, atribuído a 30 dos 78 cursos participantes na região Sul.
O conceito 2 foi atribuído a 16 cursos ( 20,5 \% ), e os conceitos 4 e 5, a 23 ( 29,5 \% ) e a seis cursos ( 7,7 \% ), respectivamente.
Dois cursos ficaram Sem Conceito ( SC ) na região Sul, e um curso recebeu o conceito 1.
Já dos 49 cursos participantes na região Centro-Oeste ( 11,3 \% do total nacional ), 22 ( 44,9 \% em termos regionais ) receberam conceito 2, o conceito modal, como já destacado.
Quatro cursos ( 8,2 \% ) obtiveram o conceito 1 ; 14 cursos ( 28,6 \% ), o conceito 3 ; e sete ( 14,3 \% ), o conceito 4.
Nessa região, dois cursos ficaram Sem Conceito ( SC ), e nenhum recebeu o conceito 5.
No Gráfico 5.1, é apresentada a distribuição cumulativa do Conceito Enade, segundo a Grande Região.
Em um gráfico de distribuição cumulativa, em cada valor de abscissa ( eixo x ), o valor apresentado no eixo das ordenadas ( eixo y ) é o valor acumulado da variável ( Conceito Enade ).
Observando -se a linha vertical pontilhada, que ajuda a visualizar o percentual de conceitos acumulados até o valor 3, ou seja, os conceitos 1, 2 e 3, pode -se notar que, por exemplo, a região Norte ( linha azul ) apresentou 100 \% dos cursos com o conceito 3 ou abaixo desse valor.
A região Sul ( linha rosa ), por outro lado, apresentou 61,8 \% dos cursos na mesma situação : com o conceito 3 ou abaixo.
Destaca -se que, em algumas regiões, o conceito 5 pode não ter sido alcançado por nenhum curso.
Foi o que ocorreu na região Norte ( linha azul ), na qual todos os cursos obtiveram conceitos menores ou iguais a 3, tendo a poligonal alcançado 100 \% no conceito 3, e na região Centro-Oeste ( linha laranja ), na qual todos os cursos obtiveram conceitos menores ou iguais a 4, tendo a poligonal alcançado 100 \% no conceito 4.
Em linhas gerais, em um gráfico de distribuição cumulativa, poligonais mais à esquerda representam as Grandes Regiões com a pior distribuição, e poligonais mais à direita, as Grandes Regiões com os melhores conceitos.
A região Norte ( linha azul ) apresentou os cursos com os conceitos mais baixos, ao passo que as regiões Sul ( linha rosa ) e Sudeste ( linha vermelha ) apresentaram a distribuição com valores maiores.
Cursos em IES situadas na região Nordeste ( linha verde ) e na região Centro-Oeste ( linha laranja ) apresentaram uma situação intermediária. 5.2 CONCEITOS POR CATEGORIA ADMINISTRATIVA, POR MODALIDADE DE ENSINO E GRANDE REGIÃO Na Tabela 5.2, é apresentada a distribuição dos cursos de Matemática ( Licenciatura ) participantes do Enade/2021, por Categoria Administrativa e Modalidade de ensino, de acordo com os conceitos por eles alcançados, segundo as Grandes Regiões brasileiras.
Dos 435 cursos participantes, 335 ( 77,0 \% ) eram ministrados em Instituições Públicas, e 100 ( 23,0 \% ), em Instituições Privadas.
Quanto à Modalidade de Ensino, 76 cursos ( 17,5 \% ) foram a Distância, e 359 ( 82,5 \% ), Presenciais.
De acordo com as informações da Tabela 5.2, em termos nacionais, entre Instituições Públicas e Privadas, observa -se que, dos 17 cursos avaliados com o conceito 5, 15 foram oferecidos em IES Públicas, e dois, em IES Privadas.
Dos 335 cursos participantes de IES Públicas, o conceito 3 foi o valor modal, atribuído a 116 cursos ( 34,6 \% ).
Dos demais cursos de Instituições Públicas participantes, 14 cursos ( 4,2 \% ) foram avaliados com o conceito 1 ; 113 cursos, com o conceito 2 ; 70 cursos, com o conceito 4 ; e, como já comentado, 15 cursos ( 4,5 \% ), com o conceito 5.
Nessa categoria, sete cursos ficaram Sem Conceito ( SC ).
Na rede privada, o conceito modal foi o 2, atribuído a 40 cursos, o correspondente a 40,0 \% dos 100 cursos da categoria.
Dos demais cursos participantes, cinco ( 5,0 \% ) receberam o conceito 1, e 32 ( 32,0 \% ), o conceito 3.
O conceito 4 foi atribuído a 13 cursos ( 13,0 \% ), e o conceito 5, a dois cursos ( 2,0 \% ).
Nessa Categoria Administrativa, oito cursos ( 8,0 \% ) ficaram Sem Conceito ( SC ).
Dos 359 cursos participantes da modalidade presencial, o conceito 3 foi o valor modal, atribuído a 130 cursos ( 36,2 \% ).
Dos demais cursos nessa Modalidade de Ensino, 13 cursos ( 3,6 \% ) receberam o conceito 1 ; 112 cursos ( 31,2 \% ) receberam o conceito 2 ; 78 cursos ( 21,7 \% ), o conceito 4 ; e 17 cursos ( 4,7 \% ) ; o conceito 5.
Nessa Modalidade de Ensino, nove cursos ( 2,5 \% ) ficaram Sem Conceito ( SC ).
Foram 76 os cursos a Distância informados, e 41 deles obtiveram o conceito modal 2.
Nessa Modalidade de Ensino, nenhum recebeu o conceito 5, e seis cursos ficaram Sem Conceito ( SC ).
Os demais receberam os conceitos 1 ( seis cursos ), 3 ( 18 cursos ) e 4 ( cinco cursos ).
Na análise por região, observa -se que, na região Norte, as Instituições Privadas participaram com três cursos ( 8,8 \% do total regional ), e um ficou Sem Conceito ( SC ).
Os demais obtiveram o conceito 2.
As Instituições Públicas participaram com 31 cursos na região Norte ( 91,2 \% do total regional ), e 20 cursos obtiveram o conceito modal 2.
Os demais receberam os conceitos 1 ( quatro cursos ) e 3 ( seis cursos ), e um curso ficou Sem Conceito ( SC ) nessa combinação de categoria e região.
Na região Norte, 28 dos 34 cursos foram oferecidos na Modalidade Presencial.
Para esses, a moda foi o conceito 2, com 19 cursos.
Os demais cursos presenciais da região receberam os conceitos 1 ( um curso ) e 3 ( seis cursos ).
Ficaram Sem Conceito ( SC ) dois cursos, e nenhum recebeu os conceitos 4 e 5.
Já para os seis cursos informados na Modalidade a Distância, a distribuição foi bimodal, com três cursos recebendo o conceito 1, e três, o conceito 2.
Na região Nordeste, oito dos 132 cursos participantes eram da rede privada, o equivalente a 6,1 \% do total da região.
Entre os cursos oferecidos por IES Privadas no Nordeste, quatro cursos foram avaliados com conceito 2, o conceito modal.
Um curso recebeu o conceito 1, dois cursos receberam o conceito 3, e um curso ficou Sem Conceito ( SC ).
Nessa combinação de categoria e região, nenhum curso recebeu os conceitos 4 e 5.
As Instituições Públicas da região Nordeste participaram com 124 cursos ( 93,9 \% do total da região ).
A distribuição nessa categoria foi bimodal, com os conceitos 2 e 3 sendo atribuídos a 53 cursos cada.
Sete cursos receberam o conceito 1 ; oito cursos, o conceito 4 ; e um curso, o conceito 5.
Dois cursos ficaram Sem Conceito ( SC ).
Na região Nordeste, a maioria dos cursos foi oferecida na Modalidade Presencial : 111 de 132 cursos.
Para esses, a moda foi o conceito 3, com 52 cursos.
Os demais cursos presenciais da região obtiveram os conceitos 1 ( sete cursos ), 2 ( 42 cursos ), 4 ( oito cursos ) e 5 ( um curso ).
Um curso ficou Sem Conceito ( SC ) nessa Modalidade na região Nordeste.
Já para os 21 cursos a Distância da região Nordeste, o conceito modal foi o 2, atribuído a 15 cursos.
Os conceitos 1 e 3 foram atribuídos a um e a três cursos, respectivamente.
Nessa combinação de Modalidade de Ensino e região, dois cursos ficaram Sem Conceito ( SC ), e nenhum recebeu os conceitos 4 e 5.
Na região Sudeste, a proporção de cursos da rede privada foi de 37,3 \%, correspondendo a 53 dos 142 cursos participantes.
Nessa categoria e região, o conceito modal foi o 2, atribuído a 26 cursos, e quatro cursos ficaram Sem Conceito ( SC ).
Os demais foram avaliados com os seguintes conceitos : 1 ( um curso ), 3 ( 14 cursos ), 4 ( sete cursos ) e 5 ( um curso ).
Nos 89 cursos oferecidos por Instituições Públicas na região Sudeste, o conceito modal foi o 4, atribuído a 38 cursos.
Nessa região e categoria, os demais conceitos atribuídos foram : 1 ( um curso ), 2 ( dez cursos ), 3 ( 29 cursos ) e 5 ( nove cursos ).
Também na região Sudeste, predominaram os cursos na Modalidade Presencial : 110 de 142 cursos.
Para esses, a moda foi o conceito 4, com 42 cursos.
Os demais cursos Presenciais da região obtiveram os conceitos 1 ( um curso ), 2 ( 19 cursos ), 3 ( 35 cursos ) e 5 ( dez cursos ).
Ficaram Sem Conceito ( SC ) três cursos.
Na região Sudeste, 32 cursos foram oferecidos na Modalidade Educação a Distância, e o conceito modal nessa Modalidade de Ensino foi o 2, atribuído a 17 cursos.
Ainda nessa Modalidade, um curso ficou com o conceito 1, oito cursos receberam o conceito 3, e três cursos, o conceito 4, não tendo nenhum curso recebido o conceito 5.
Na Modalidade a Distância, ficaram Sem Conceito ( SC ) três cursos.
Na região Sul, as Instituições Privadas participaram com 29 dos 78 cursos ( 37,2 \% do total regional ).
Desses, 14 cursos obtiveram conceito 3, o conceito modal.
Nessa combinação de Categoria Administrativa e Grande Região, um curso recebeu o conceito 1 ; um curso recebeu o conceito 5 ; seis cursos receberam o conceito 2 ; e outros seis, o conceito 4.
Um curso ficou Sem Conceito ( SC ).
As Instituições Públicas na região Sul participaram com 49 cursos ( 62,8 \% dos cursos da região ), e 17 cursos receberam o conceito modal 4.
Aos demais cursos foram atribuídos os conceitos 2 ( dez cursos ), 3 ( 16 cursos ) e 5 ( cinco cursos ).
Um curso da região Sul oferecido por IES Pública ficou Sem Conceito ( SC ), e nenhum recebeu o conceito 1.
Dos 78 cursos participantes da região Sul, 65 eram presenciais.
O conceito modal para essa Modalidade de Ensino foi o 3, atribuído a 23 cursos.
O conceito 1 foi atribuído a um curso ; o conceito 2, a 12 cursos ; o conceito 4, a 21 cursos ; e o conceito 5, a seis cursos.
Dois cursos Presenciais da região Sul ficaram Sem Conceito ( SC ).
Já dos 13 cursos oferecidos nessa região na Modalidade a Distância, sete receberam o conceito modal 3, quatro receberam o conceito 2, e dois, o conceito 4.
Na região Centro-Oeste, sete dos 49 cursos participantes eram de Instituições Privadas ( 14,3 \% em termos regionais ).
A distribuição, nessa região e categoria, foi multimodal, com os conceitos 1, 2 e 3 sendo atribuídos a dois cursos cada ; o curso restante ficou Sem Conceito ( SC ).
Quanto aos 42 cursos oferecidos por Instituições Públicas na região Centro-Oeste ( 85,7 \% do total regional ), 20 cursos apresentaram o conceito 2, conceito modal ; já os demais cursos receberam os conceitos 1, 3 e 4, atribuídos, respectivamente a, dois, 12 e sete cursos.
Um curso ficou Sem Conceito ( SC ), e nenhum recebeu o conceito 5 nessa combinação de categoria e região.
Na região Centro-Oeste, participaram 45 cursos presenciais, e o conceito modal foi o 2, recebido por 20 cursos.
Os demais receberam os conceitos 1 ( três cursos ), 3 ( 14 cursos ) e 4 ( sete cursos ).
Um curso ficou Sem Conceito ( SC ), e nenhum recebeu o conceito 5 nessa combinação de Modalidade de Ensino e região.
Já dos quatro cursos a Distância da região Centro-Oeste, dois receberam o conceito modal 2, um recebeu o conceito 1, e um curso ficou Sem Conceito ( SC ).
No Gráfico 5.2, apresenta -se a distribuição do Conceito Enade, segundo a Categoria Administrativa da IES.
Os cursos em IES Públicas ( linha azul ) apresentam uma poligonal mais à direita do que os em IES Privadas ( linha verde ) e, consequentemente, uma distribuição de Conceito Enade com valores maiores. 5.3 CONCEITOS POR ORGANIZAÇÃO ACADÊMICA E GRANDE REGIÃO Na Tabela 5.3, encontra -se a distribuição dos conceitos atribuídos aos cursos participantes do Enade/2021, na Área de Matemática ( Licenciatura ), por Organização Acadêmica, segundo as Grandes Regiões brasileiras.
Dos 435 cursos de Matemática ( Licenciatura ) participantes, 282 eram oferecidos em Universidades ; 26, em Centros Universitários ; 35, em Faculdades ; e 92, em CEFET/IF.
Esta distribuição corresponde a, respectivamente, 64,8 \%, 6,0 \%, 8,0 \% e 21,1 \% dos cursos.
De acordo com os dados apresentados, de todos os 17 cursos avaliados com o conceito 5, 15 eram vinculados a Universidades.
Os cursos vinculados a esse tipo de Organização Acadêmica tiveram o conceito 3 como conceito modal, atribuído a 100 cursos ( 35,5 \% ).
Os demais cursos vinculados a Universidades receberam os conceitos 1 ( oito cursos ), 2 ( 97 cursos ), 4 ( 55 cursos ) e 5 ( 15 cursos, como já mencionado ).
Sete cursos ficaram Sem Conceito ( SC ).
Nos 26 cursos em Centros Universitários, o conceito modal foi o 2, atribuído a 12 cursos ( 46,2 \% ).
Cinco cursos vinculados a esse tipo de Organização Acadêmica ficaram Sem Conceito ( SC ), e os outros cursos receberam os conceitos 1 ( um curso ), 3 ( seis cursos ) e 4 ( dois cursos ).
Nenhum curso nesse tipo de Organização Acadêmica recebeu o conceito 5.
Dos 35 cursos mantidos por Faculdades, dez ( 28,6 \% ) receberam conceito 2, o conceito modal.
O segundo conceito mais frequentemente atribuído aos cursos oferecidos por essa categoria de IES foi o conceito 1, recebido por nove cursos.
Na sequência, sete cursos obtiveram conceito 3 ; cinco cursos, o conceito 4, e um curso recebeu o conceito 5.
Dos cursos oferecidos em Faculdades, três ficaram Sem Conceito ( SC ).
Foram 92 os cursos mantidos por CEFET/IF, e o conceito modal para esse tipo de Organização Acadêmica foi o 3 ( 35 cursos ).
Nenhum curso ficou Sem Conceito ( SC ).
Os conceitos 1 e 5 foram atribuídos a um curso cada, e os conceitos 2 e 4 a, respectivamente, 34 cursos e 21 cursos.
Considerando -se, separadamente, as regiões brasileiras, verifica -se que, na região Norte, as Universidades participaram com 21 dos 34 cursos avaliados, e foi atribuído o Considerando -se, separadamente, as regiões brasileiras, verifica -se que, na região Norte, as Universidades participaram com 21 dos 34 cursos avaliados, e foi atribuído o conceito modal 2 a 12 cursos.
Os demais conceitos foram 1 ( três cursos ) e 3 ( cinco cursos ).
Um curso ficou Sem Conceito ( SC ), e nenhum curso recebeu os conceitos 4 e 5 nessa combinação de categoria e região.
Os Centros Universitários da região Norte foram representados por um curso, o qual obteve o conceito 2.
As Faculdades participaram com dois cursos na região Norte, dos quais um ficou Sem Conceito ( SC ), e um recebeu o conceito 2.
Os CEFET/IF tiveram dez cursos avaliados na região Norte, dos quais oito obtiveram o conceito modal 2, e os dois restantes, os conceitos 1 e 3.
Na região Nordeste, as Universidades participaram com 85 dos 132 cursos da Área de Matemática ( Licenciatura ) oferecidos.
Dois desses 85 cursos ficaram Sem Conceito ( SC ).
O conceito modal 3 foi atribuído a 38 cursos.
Dois cursos obtiveram o conceito 1 ; 35 cursos o conceito 2 ; sete cursos, o conceito 4 ; e um curso obteve o conceito 5.
Os Centros Universitários contaram com quatro cursos participantes na região Nordeste, dos quais dois receberam o conceito modal 2, um recebeu o conceito 3, e um ficou Sem Conceito ( SC ).
As Faculdades foram representadas por nove cursos na região Nordeste, seis dos quais receberam conceito 1, o conceito modal.
Os três cursos restantes nessa Organização Acadêmica receberam o conceito 2.
Foram 34 os cursos participantes mantidos por CEFET/IF na região Nordeste, e o conceito modal 2 foi atribuído a 17 cursos.
Os conceitos 3 e 4 foram atribuídos a, respectivamente 16 cursos e um curso.
Nenhum curso ficou Sem Conceito ( SC ) ou recebeu os conceitos 1 ou 5.
Na região Sudeste, as Universidades concentraram 83 dos 142 cursos de Matemática ( Licenciatura ) da região.
Dos cursos oferecidos em Universidades, nessa região, o conceito modal foi o 3, atribuído a 26 cursos.
Os demais cursos receberam os conceitos 1 ( um curso ), 2 ( 21 cursos ), 4 ( 25 cursos ) e 5 ( oito cursos ).
Dois cursos ficaram Sem Conceito ( SC ).
Os Centros Universitários participaram com 14 cursos na região Sudeste, dos quais sete obtiveram o conceito modal 2.
O conceito 3 foi atribuído a três cursos, e o conceito 4, a um curso.
Três cursos ficaram Sem Conceito ( SC ), e nenhum recebeu os conceitos 1 ou 5 nessa combinação de Organização Acadêmica e região.
Já as Faculdades foram representadas por 16 cursos na região Sudeste, e a distribuição foi bimodal, com os conceitos 2 e 3 sendo atribuídos a quatro cursos cada.
Os demais cursos receberam os conceitos 1 ( um curso ), 4 ( cinco cursos ) e 5 ( um curso ), tendo um curso ficado Sem Conceito ( SC ).
Dos 29 cursos avaliados de CEFET/IF na região Sudeste, 14 receberam o conceito 4, valor modal.
Quatro cursos receberam o conceito 2, dez cursos receberam o conceito 3, e um curso, o conceito 5.
Nenhum curso ficou Sem Conceito ( SC ) ou recebeu o conceito 1 nessa combinação de organização acadêmica e região.
Dos 78 cursos da região Sul, 56 eram mantidos por Universidades, para os quais o conceito modal foi o 3, atribuído a 23 cursos.
Os demais cursos receberam os conceitos 2 ( dez cursos ), 4 ( 16 cursos ) e 5 ( seis cursos ).
Nesse tipo de organização, um curso da região Sul ficou Sem Conceito ( SC ), e nenhum recebeu o conceito 1.
Dos cinco cursos participantes de Centros Universitários da região Sul, um ficou Sem Conceito ( SC ).
Os conceitos restantes foram atribuídos a dois cursos ( conceito modal 2 ), um curso ( conceito 3 ) e um curso ( conceito 4 ).
Dos cinco cursos oferecidos em Faculdades na região Sul, três receberam o conceito modal 3.
Os conceitos 1 e 2 tiveram a mesma quantidade de cursos, um cada.
Nessa região, nenhum curso oferecido em Faculdades ficou Sem Conceito ( SC ) ou obteve os conceitos 4 e 5.
Os CEFET/IF participaram com 12 cursos, dos quais seis receberam o conceito modal 4, e os demais receberam os conceitos 2 e 3 ( três cursos cada ).
Na região Centro-Oeste, 37 dos 49 cursos oferecidos eram mantidos por Universidades.
Para esse tipo de Organização Acadêmica, um curso ficou Sem Conceito ( SC ), e o conceito modal 2 foi obtido por 19 cursos.
Os demais conceitos atribuídos foram : 1 ( dois cursos ), 3 ( oito cursos ) e 4 ( sete cursos ).
Nessa região, nenhum curso recebeu o conceito 5.
Os Centros Universitários da região Centro-Oeste contaram com dois cursos, que receberam os conceitos 1 e 3.
Dos três cursos oferecidos em Faculdades na região Centro-Oeste, um ficou Sem Conceito ( SC ), e os demais receberam os conceitos 1 e 2.
Já dos sete cursos oferecidos em CEFET/IF na região Centro-Oeste, cinco receberam o conceito modal 3, e dois, o conceito 2.
No Gráfico 5.3, apresenta -se a distribuição do Conceito Enade dos cursos de Matemática ( Licenciatura ), segundo a Organização Acadêmica da IES.
Os cursos em Universidades ( linha azul ) e em CEFET/IF ( linha laranja ) apresentam uma poligonal mais à direita do que as demais, denotando uma melhor distribuição de Conceitos Enade.
No outro extremo, encontram -se os cursos mantidos por Centros Universitários ( linha verde ), apresentando a respectiva poligonal mais à esquerda e, consequentemente, uma distribuição de Conceitos Enade com menores valores.
Os cursos mantidos por Faculdades ( linha vermelha ) apresentam a poligonal mais central e, consequentemente, uma distribuição de Conceitos Enade com valores intermediários.

% matched lemmas: 
\end{textsample}
